\chapter{横断面研究与队列研究的设计}
\chapterauthors{Thomas B. Newman、Warren S. Browner、Steven R. Cummings}
\sourcepage{217}
观察性研究主要有两个目的：一是描述性目的，考察人群中变量的分布；二是分析性目的，考察变量之间的关联。本章介绍两种基本的观察性设计，按测量的时间安排加以区分。

在\term{横断面研究}中，研究者在同一时点或较短时间内完成全部测量。他从人群中抽取样本，考察样本内各变量的分布，有时根据生物学合理性、既往资料和研究假设，将变量指定为预测变量或结局变量。例如，要研究体重与血压的关系，可以在每位参与者的一次门诊访视中测量这两个变量，考察体重较高者是否更可能患高血压。

在\term{队列研究}中，研究开始时确定一组参与者（即“队列”），并在一段时间内进行测量。其界定性特征是：在一开始组建这组人，并进行纵向（随时间）随访。正如横断面研究的预测变量与结局变量在结构上没有区别，队列研究中，特别关注的预测变量（常为某种暴露或治疗）与其他结局预测变量（常称为\term{协变量}，covariates）在结构上也没有区别。例如，可以测量一组参与者的鸡蛋摄入量及协变量（如其他饮食习惯、年龄和吸烟情况），然后随访，确定鸡蛋摄入与心血管疾病发生的关系；同一研究也能量化各协变量的作用（1）。本章讨论前瞻性、回顾性以及多队列设计，并介绍如何尽可能保留随访对象及队列研究统计分析的基本方法。

\section{横断面研究}
\subsection{基本设计}
横断面研究在大致同一时间完成所有测量，没有随访期（图8.1）。它很适合描述变量及其分布模式。例如，美国第一次全国健康与营养检查调查（National Health and Nutrition Examination Survey，NHANES-I）于1971—1975年访谈并检查了一个旨在代表美国全部1～74岁人口的样本。这项横断面研究是了解当时美国人口健康和生活习惯的重要来源，给出了不同人口学群体吸烟率等估计。此后开展了NHANES-II（1976—1980）、西班牙裔HANES（1982—1984）和NHANES-III（1988—1994）；自1999年起，NHANES持续进行，每两年修订一次（2）。\sourcepage{218}美国疾病控制与预防中心（Centers for Disease Control，CDC；\url{www.cdc.gov/nchs/nhanes.htm}）提供全部NHANES数据集和资料收集表，供公众使用。

\providecommand{\chEightCross}{%
\begin{tikzpicture}[font=\small,>=Stealth]
\draw[dashed,dcrred] (0,0) ellipse (3.0 and 2.8);
\node[text=dcrred] at (0,3.05) {可接触人群};
\node[align=center] at (0,1.95) {符合纳入标准的样本};
\draw[->] (1.7,1.9) to[bend left] (2,1.3);
\draw[dcrred] (-2.7,-1.3) rectangle (2.7,1.3);
\node[text=dcrpurple,align=center] at (0,.7) {测量当前变量\\获取既往资料（如相关）};
\node[draw,align=center] (x) at (-1.35,-.55) {符合\\排除标准};
\node[draw,align=center] (y) at (1.35,-.55) {不同意\\参加};
\node (e) at (0,-2) {排除};\draw[->] (e)--(x);\draw[->] (e)--(y);
\end{tikzpicture}}
\providecommand{\chEightCohort}[1]{%
\begin{tikzpicture}[font=\footnotesize,>=Stealth,x=.9cm,y=.9cm]
\draw[dashed,dcrred] (0,0) ellipse (2.4 and 2.2);
\node[text=dcrred] at (0,2.5) {可接触人群};
\node[align=center] at (0,1.6) {符合选择标准的样本};
\draw[->] (1.1,1.6) to[bend left] (1.5,1.05);
\draw[dcrred] (-2.1,-1.05) rectangle (2.1,1.05);
\node[text=dcrpurple,align=center] at (0,0) {测量基线变量\\[9pt]\ifnum#1=0 获取既往资料（如相关）\else 保存标本（可选）\fi};
\draw[thick,dcrteal,->] (2.1,.8)--(9.1,.8);
\node[text=dcrteal,align=center] at (5.6,1.15) {随访参与者};
\draw[thick,dcrteal,->] (2.1,-.8)--(9.1,-.8);
\ifnum#1>0
\ifnum#1=2
\foreach \x in {3.05,5.45,7.7}{\node[text=dcrgold,align=center,text width=9mm,inner sep=1pt,font=\fontsize{8}{9.5}\selectfont] at (\x,0) {结局\\发生时\\测量};}
\else
\node[text=dcrgold,align=center] at (5.7,0) {结局发生时测量};
\fi
\foreach \x in {3.2,5.8,8}{\draw[->,dcrteal] (\x,-.8) to[bend left=30] (\x+.6,-2.05);}
\node[text=dcrteal] at (6,-2.4) {失访};
\else
\draw[->,dcrteal] (6,-.8) to[bend left=35] (7,-1.8);
\node[text=dcrteal] at (6,-2.3) {极少失访或无人失访};
\fi
\ifnum#1=2
\foreach \x in {4.1,6.8}{\node[draw=dcrpurple,fill=white,text=dcrpurple,rotate=90] at (\x,0) {更新预测因素};}
\fi
\node[draw=dcrgold,text=dcrgold,align=center] at (9.6,0) {\ifnum#1=0 随访结束时\\测量结局\else 可选择在随访\\结束时测量结局\fi};
\draw[dcrteal,thick,->] (2.1,-3)--(9,-3);
\node[fill=white,text=dcrteal] at (5.5,-3) {随访时间\ifnum#1=0 （短）\fi};
\draw[dcrteal] (-2.1,-1.05)--(-2.1,-2.7)--(2.1,-2.7)--(2.1,-1.05);
\node[text=dcrteal] at (0,-3) {研究起始};
\end{tikzpicture}}
\providecommand{\chEightDouble}{%
\begin{tikzpicture}[font=\footnotesize,>=Stealth,x=.9cm,y=.9cm]
\foreach \yy in {0,-5.5}{
\begin{scope}[yshift=\yy cm]
\draw[dashed,dcrred] (0,0) ellipse (2.4 and 2.1);
\node[align=center] at (0,1.6) {符合选择标准的样本};
\draw[->] (1.1,1.6) to[bend left] (1.5,1);
\draw[dcrred] (-2.1,-1) rectangle (2.1,1);
\node[text=dcrpurple,align=center] at (0,0) {测量基线变量\\保存标本（很少）};
\draw[thick,dcrteal,->] (2.1,.7)--(9,.7);
\draw[thick,dcrteal,->] (2.1,-.7)--(9,-.7);
\node[text=dcrteal] at (5.4,1) {随访参与者};
\node[text=dcrgold] at (5.6,.1) {结局发生时测量};
\foreach \x in {3.3,6,8}{\draw[->,dcrteal] (\x,-.7) to[bend left=30] (\x+.6,-1.9);}
\node[text=dcrteal] at (6,-2.3) {失访};
\node[draw=dcrgold,text=dcrgold,align=center] at (9.7,0) {可选择在随访\\结束时测量结局};
\end{scope}}
\node[text=dcrred,align=center] at (0,2.65) {可接触的暴露人群\\（人数常有限）};
\node[text=dcrred,align=center] at (0,-2.8) {可接触的未暴露或不同暴露人群\\（通常较大）};
\node[fill=white,align=center] at (0,-4.25) {符合选择标准的样本\\（人数比可为1:1、2:1等）};
\draw[dcrteal,thick,->] (2.1,-8.8)--(9,-8.8);
\node[fill=white,text=dcrteal] at (5.4,-8.8) {随访时间};
\node[text=dcrteal] at (0,-8.8) {研究起始};
\end{tikzpicture}}

\begin{figure}[htbp]\centering
\chEightCross
\caption*{\textbf{图8.1\quad 横断面研究。}步骤为：从可接触人群中抽取符合纳入标准的参与者；排除符合排除标准或不愿同意参加者；测量预测变量与结局变量的当前值，通常辅以既往资料。}
\end{figure}

横断面研究可以考察关联，但将哪些变量称为预测变量、哪些称为结局变量，取决于研究者的因果假设，而非研究设计。年龄、种族和出生时指定性别等固有因素不会被其他变量改变，因此总是预测变量，这种选择很容易。其他变量则可能有两种解释。例如，NHANES-III发现儿童肥胖与看电视的时长存在横断面关联（3）；将哪一个称为预测变量、另一个称为结局，取决于研究者的因果假设。

队列研究有纵向时间维度，可以估计\term{发病}（incidence，即一段时间内发生某种疾病或状况的人所占比例）；横断面研究则只能提供\term{患病率}（prevalence，即某时点有该疾病或状况的人所占比例）的信息（表8.1）。临床人群的横断面研究（如急诊腹痛患者）能帮助临床医师估计眼前患者患某种疾病（如阑尾炎）的可能性；患病率越高，疾病的“先验概率”越高（第12章）。横断面研究也能帮助卫生规划者了解有多少人患某些疾病，以便配置足够的照护资源。分析时可以比较有、无某种暴露者的结局患病率，得到\term{患病率比}（prevalence ratio），即横断面研究中与风险比相应的指标。

\begin{table}[htbp]\centering\small
\caption*{\textbf{表8.1\quad 发病与患病率的区别}}
\begin{tabularx}{\textwidth}{L{22mm}L{36mm}Y}\toprule
研究类型&统计量&定义\\\midrule
横断面&患病率&某时点\textbf{患有}某疾病或状况的人数／处于风险中的人数\\
队列&累积发病率或发病比例&给定时间内\textbf{发生或曾发生}某疾病的人数／处于风险中的人数\\
队列&发病密度&\textbf{发生或曾发生}某疾病的人数／（处于风险中的人数$\times$处于风险中的时间）\\\bottomrule
\end{tabularx}
\end{table}
\sourcepage{219}
横断面研究有时描述曾经有某种暴露，或曾经患某疾病、出现某状况的比例。不过，年龄越大，经历暴露或患病的机会就越多。因此，考察这类“曾经有过”的变量之间的关联时，必须考虑年龄，或其他与暴露机会较多相关的因素，例如职业暴露研究中的从业年限。例8.1在接触不同程度演员吸烟电影的儿童中考察曾尝试吸烟的比例。看过更多电影的儿童通常年龄也更大，因而有更多机会尝试吸烟，所以多变量分析中调整年龄十分重要（见第10章）。

\begin{bookbox}{例8.1\quad 横断面研究}
Sargent等（4）希望确定，观看演员吸烟的电影是否与开始吸烟有关。步骤如下：
\begin{enumerate}
\item \textbf{界定选择标准并招募人群样本。}研究者通过随机数字拨号调查了6522名10～14岁的美国儿童。
\item \textbf{测量预测变量和结局变量。}量化532部热门电影中的吸烟内容；对每位参与者，从中随机抽取50部，询问他看过哪些。还询问年龄、种族、性别、父母吸烟和教育情况、寻求刺激的倾向（如“我喜欢做危险的事”）和自尊（如“我希望自己是另外一个人”）等协变量。结局是儿童是否曾尝试吸过一支烟。
\end{enumerate}
曾尝试吸烟的比例，从电影吸烟暴露最低四分位组的2\%，上升到最高四分位组的22\%。调整年龄和其他混杂因素后，差异仍有统计学意义；作者估计，38\%的吸烟开始可归因于观看演员吸烟的电影。此后多个国家的研究证实，电影中的吸烟内容是青少年开始吸烟的重要促成因素（5—7）。
\end{bookbox}

\subsection{横断面研究的优点与缺点}
横断面研究的一大优点是无需等待结局发生，因此速度快、费用低，并避免了失访问题。\sourcepage{220}另一优点是，几乎不增加成本，就能将其作为队列研究或临床试验的第一步。其结果界定研究组的基线人口学和临床特征，使后续纵向研究能够排除已经出现结局者，有时还能发现有意义的横断面关联。横断面研究也适用于诊断试验研究：患者已经患有关注的疾病，重点是检查诊断该疾病的准确性（第13章）。

对于在短期内相互影响的特征，横断面设计是考察其关联的最佳方法。例如，睡眠模式与认知测验之间的关联，最好使用同一天的测量研究（8）。

但如前所述，横断面资料往往难以确定因果关系。它只测量患病率，而非发病，因此推断疾病的病因、预后或自然史时必须谨慎。与患病率有关的因素可能是病因，也可能只是通过影响死亡或疾病消退的可能性，与病程长短相关。\footnote{实际上，疾病患病率是发病率与（平均）病程的乘积：$P=I\times D$。}例如，慢性肾病患病率不仅取决于发病，还取决于患病后的生存。鉴于肥胖与透析患者较好的生存相关（9），对现患慢性肾衰竭预测因素的横断面研究，会高估肥胖与慢性肾衰竭发生之间的关联。

\subsection{连续调查}
研究者有时会在同一人群中反复开展横断面调查，例如每5年一次。这种设计可以推断模式随时间的变化，前述NHANES就是如此。有时，它能发现较短时间内患病率的显著变化。例如，McMillen等（10）每年对有代表性的美国成人样本开展横断面调查，发现2010—2013年，年轻成人（18～24岁）使用电子烟的比例从0\%增至14.2\%。

连续调查具有纵向时间框架，但不是队列研究，因为每次都重新抽取样本。因此不能评估个体内的变化；出生、死亡和迁移使人们进入或离开人群（以及样本），也可能影响结果。

\section{队列研究}
“Cohort”在古罗马指共同进军的一队士兵。在临床研究中，队列指一开始就明确界定、并随时间随访的一组参与者。

\subsection{前瞻性与回顾性研究}
队列研究常分为\term{前瞻性}（prospective）与\term{回顾性}（retrospective，字面意思分别为向前看和向后看）。但这些词的定义不一（11），两者之间也没有明确的二分界线（例8.2、8.3）：同一队列研究可以对一些变量向前看，对另一些向后看，即“双向性队列研究”（ambidirectional cohort studies）（12）。因此，这些词最好用于描述研究的某些方面，而非给整项研究贴上非此即彼的标签。具有较多前瞻性特征的研究较少受偏倚影响，但往往更困难、费用更高，因为研究者必须组建队列并随访，有时持续多年。\sourcepage{221}新研究者可能难以实施，不过，常可使用正在进行或已经完成的队列研究中的资料或标本，利用前瞻性研究的优势（第16章）。

\begin{bookbox}{例8.2\quad 一项（主要为）前瞻性队列研究}
经典的护士健康研究（Nurses’ Health Studies，NHS）自1976年原始研究开始招募起，一直考察女性常见病的流行病学。（2015年起，NHS-3也招募男护士。）（13）

2020年，Wang等报告了护士健康研究II（NHS-II）中月经周期规律性与70岁前死亡风险的关联。步骤如下：
\begin{enumerate}
\item \textbf{界定选择标准并组建队列。}NHS-II于1989年从美国14个州招募116429名25～42岁女性注册护士。
\item \textbf{测量预测变量，包括潜在混杂因素。}NHS-II特别关注口服避孕药和其他生殖相关疾病危险因素。1989年基线时，要求参与者回忆14～17岁和18～22岁时的月经周期特征。1993年，这些当时29～46岁的女性报告其当前周期；1993年的其他协变量值被作为基线测量。
\item \textbf{随访队列并测量结局。}通过州生命统计记录、定期检索全国死亡索引，或近亲及邮政机构报告，确定死亡。
\end{enumerate}
研究发现，与通常周期为26～31天的女性相比，报告通常周期至少40天者更可能在70岁前死亡：18～22岁时周期长者的调整后风险率比为1.34（95\%置信区间1.06～1.69）；29～46岁时为1.40（95\%置信区间1.17～1.68）。调整潜在混杂因素未改变结果。

\textbf{点评：}14～17岁及18～22岁的周期特征来自招募时的回忆，在这一意义上是回顾性的；当前月经周期的测量则是前瞻性的。结局为70岁前全因死亡，其测量的偏倚风险很低，而且结局发生不可能影响暴露测量。
\end{bookbox}

\begin{bookbox}{例8.3\quad 一项（主要为）回顾性队列研究}
若在症状出现数小时内给予，抗栓（“稀释血液”）治疗可改善急性缺血性卒中的结局。Man等（14）考察了Medicare患者“入院至用药时间”（door-to-needle time，即从到达医院到给予组织型纤溶酶原激活剂〔tissue plasminogen activator，tPA〕治疗的时间）与1年死亡的关联。步骤如下：
\sourcepage{222}
\begin{enumerate}
\item \textbf{寻找合适的既有队列。}研究纳入61426名至少65岁的Medicare受益人；他们就诊于参加Get With the Guidelines（GWTG）-Stroke质量改进项目的医院，并在症状出现4.5小时内接受tPA。
\item \textbf{收集暴露资料和协变量。}作为GWTG-Stroke项目的一部分，经过培训的医院人员收集入院至用药时间及其他临床结局预测因素；作者从美国医院协会数据库获得医院层面特征。
\item \textbf{确定结局。}通过一系列间接标识信息匹配，包括入院和出院日期、医院、患者性别和出生日期，将GWTG-Stroke记录与Medicare理赔档案连接，获得随后一年内死亡与住院资料。
\end{enumerate}
研究者以Cox比例风险模型调整基线协变量后，发现入院至用药时间与全因死亡存在剂量反应关系：每延长15分钟，调整后风险率比为1.04（95\%置信区间1.02～1.05）。（Cox模型是对数模型，因此延长1小时的估计作用，是该小时内每15分钟均乘以1.04，即$1.04^4=1.17$。）不过，较短的入院至用药时间也与多种已测量的患者和医院基线特征相关。

\textbf{点评：}因此，作出因果推断依赖于以下假定：统计模型充分调整了上述全部差异，而且与死亡相关的未测量特征不存在重要差异（第10章）。主要暴露变量入院至用药时间，是GWTG-Stroke质量改进项目中仔细测量的。假如测量者是研究团队，而非医院质量改进人员，和／或患者曾为参与研究提供同意，研究者很可能会称其为前瞻性而非回顾性队列研究。
\end{bookbox}

最纯粹的前瞻性队列研究，先确定一组有可能发生所研究结局的参与者（例如排除已经有结局者或已没有该结局发生器官者），在基线界定并测量预测因素，再按严格的结局确定方案随访，观察结局的新发生或变化。例如，在基线无周围神经病变的3000名参与者中测量血镉及其他预测因素，随后定期测量足部触觉、痛觉和振动觉，以发现周围神经病变。

某些前瞻性特征能提高研究效度（表8.2），其中最重要的是确保结局（如周围神经病变）的发生和测量不能影响预测因素（如镉水平）的测量。此外，前瞻性测量使研究者更能控制预测变量与结局变量的测量质量。例如，镉最好用特殊无金属试管收集的全血测量，而回顾性研究可能只保留了普通试管中的血清或血浆。同样，前瞻性研究可由神经科医师每半年使用标准化工具检查一次，持续10年，以确定是否发生周围神经病变；\sourcepage{223}回顾性研究可能只能依赖电子病历中的诊断。前瞻性研究质量更高，但成本也高得多，完成所需时间更长。

\begin{longtable}{L{48mm}L{101mm}}
\caption*{\textbf{表8.2\quad 使队列研究更具前瞻性的特征及其优缺点}}\\\toprule
特征&优点（$+$）与缺点（$-$）\\\midrule\endfirsthead
\toprule 特征&优点（$+$）与缺点（$-$）\\\midrule\endhead
研究者将合格参与者招入队列&$+$ 明确入选标准，并能控制招募资格。\newline $-$ 招募、资料与标本收集等费用。\\
研究者在基线检查结局，排除已出现结局者&$+$ 避免结局先于暴露，甚至引起暴露。\newline $-$ 在基线进行终点的体格测量可能昂贵。\newline $-$ 为招到足够合格者，必须筛查更多人，增加时间和费用。\\
研究者测量暴露与协变量，不依赖既往测量&$+$ 尽可能提高测量的准确度和精密度。\newline $+$ 尽可能减少暴露与协变量的缺失资料。\newline $-$ 进行这些测量所需时间和费用。\\
在结局发生前测量暴露&$+$ 防止结局影响暴露的测量。\\
结局发生时或随访结束时由研究者测量，不依赖已完成的结局测量&$+$ 防止结局影响暴露的测量。\newline $+$ 更能控制结局的确定。\newline $+$ 可能减少结局缺失资料。\newline $-$ 耗时（必须等待结局发生）且昂贵。\\\bottomrule
\end{longtable}

情况还更复杂：几乎所有前瞻性研究，也会测量研究开始前已开始或已发生的预测因素，例如询问既往暴饮史。虽然这一预测因素是回顾性（向后看）测量的，但测量在研究开始时进行，因而结局（如神经病变）的知晓或发生，都不可能影响该因素的发生或测量。

此外，前瞻性研究资料常被用于回答研究开始时未曾考虑的问题。例如，在骨质疏松性骨折研究中，我们考察低骨密度是否与卒中风险增加相关（15）。骨密度虽在基线按严格方案测量，但卒中原本并未被判定为研究结局。后来只在报告曾因卒中住院的女性中取得病历并审查，以验证卒中。这种确定卒中的方法，不及对全部参与者定期神经系统评估、并在任何突发神经症状时按标准程序评估卒中的方法严格。

甚至通常被视为回顾性的病例对照研究（第9章），即结局在研究开始前已经发生的研究，也可能具有前瞻性特征。例如，研究者可以制定严格的病例定义，在符合定义的病例发生时入组。\sourcepage{224}预测因素（如既往驾驶记录中的因素）也可能在结局发生前已经测量。这样可防止结局影响预测因素或其测量，提高效度。

判断研究的“前瞻性程度”时，研究何时\emph{开始}并非总是明确。重要的是主要预测因素和结局何时发生，以及研究者是否、何时开始严格测量。例如，正在进行的队列研究可能在第二次随访加入平衡测验。此后有人可能考察基线自报失眠是否影响平衡。不过，若收集基线资料时并不特别关注失眠，它可能未被严格测量；此外，缺乏基线平衡资料，也就无法确定失眠是否先于平衡困难。

\subsection{队列研究复杂程度的范围}
队列研究可从简单到复杂（表8.3）。最简单者只在基线测量预测因素，随访很短，无须担心失访，并在研究结束时，比较不同预测变量值者的结局\term{累积发病率}（cumulative incidence，即风险）（图8.2）。例如，新冠疫情最初数月内一项考察患者能否存活出院的预测因素研究（16），全部参与者结局均已知，病程短，适合这种设计。

\begingroup\footnotesize\setlength{\tabcolsep}{3pt}
\begin{longtable}{L{15mm}L{29mm}L{20mm}L{20mm}L{18mm}L{18mm}L{21mm}}
\caption*{\textbf{表8.3\quad 简单与较复杂队列研究的特征}}\\\toprule
复杂程度&暴露与协变量测量&结局测量&失访处理&发病指标&关联指标&多变量分析\\\midrule\endfirsthead
\toprule 复杂程度&暴露与协变量测量&结局测量&失访处理&发病指标&关联指标&多变量分析\\\midrule\endhead
简单&仅基线&研究结束时&必须极少，研究才有效；敏感性分析&累积发病率、发病比例（如每100名风险者的病例数）&风险比、比值比&logistic回归\\
中等&仅基线&发生时或定期确定&失访时删失；敏感性分析&发病密度（如每100风险人年的病例数）&发病密度比或风险率比&Poisson或Cox回归\\
复杂&基线并定期更新；暴露及协变量值可随时间改变，但协变量不受所研究暴露影响&发生时或定期确定&失访时删失；敏感性分析&发病密度（如每100风险人年的病例数）&风险率比&含时间依赖协变量的Cox模型\\
\sourcepage{225}最复杂&基线并定期更新；暴露及协变量值可随时间改变，且协变量可受所研究暴露影响&发生时或定期确定&失访时删失；敏感性分析&发病密度（如每100风险人年的病例数）&风险率比&边际结构模型等高级（且困难）方法。应考虑更简单的设计！\\\bottomrule
\end{longtable}\endgroup

\begin{figure}[htbp]\centering\chEightCohort{0}
\caption*{\textbf{图8.2\quad 简单队列研究。}最初步骤与横断面研究（图8.1）相似；然后随访全部参与者一段预定时间，并在期末确定谁发生结局（或重复测量结局变量）。}
\end{figure}
\sourcepage{226}
\begin{figure}[htbp]\centering\chEightCohort{1}
\caption*{\textbf{图8.3\quad 中等复杂程度的队列研究。}最初步骤类似简单队列研究（图8.2）；不同之处是随访每个人直到发生结局、失访或研究结束，并记录各事件发生的时间。}
\end{figure}

复杂程度的下一层次，是考虑随访时间，在结局发生时测量，但仍只在入队列时测量暴露（图8.3）。这样可处理随访时间不同和失访，这在持续数月以上或参与者不易联系的研究中很常见。它不仅考虑处于风险中的人数，也考虑他们处于风险中的时间，两者相乘即\term{风险人时}（person-time at risk）。分析比较的是以风险人时为分母的\term{率}（rate），而非仅以风险人数为分母的\term{风险}（risk；见“队列研究的分析方法”）。有些参与者会失访，或研究结束时尚未发生结局，因此必须处理\term{删失}（censoring），即不能完整知晓结局的情况。

若定期更新预测因素的测量（图8.4），分析更复杂，也更不直观，因为已不能简单比较暴露者与未暴露者：同一人在不同时段，可贡献不同暴露水平的人时。例如，一个参与者可能在研究中途开始做填字游戏（关注的暴露）。因此，分析通常比较暴露与未暴露\emph{人时}中的事件风险，而非暴露与未暴露\emph{人群}。若暴露效应有延迟，既往暴露史也很重要，问题尤其困难。再者，若结局发生时即测量，而预测因素仅定期更新，结局发生时的预测因素值可能已不同于最近一次测量值。护士健康研究中一个令人遗憾的例子清楚展示了这一点（例8.4）。
\sourcepage{227}
\begin{figure}[htbp]\centering\chEightCohort{2}
\caption*{\textbf{图8.4\quad 复杂队列研究。}步骤类似中等复杂程度研究（图8.3），但研究者重复测量以更新预测变量，并可在分析中同时考虑预测变量最近值和既往值。}
\end{figure}

\begin{bookbox}{例8.4\quad 研究现用者并过分追求“前瞻性”所付出的代价}
20世纪流行病学的一大失败，是提倡绝经后激素治疗以延长寿命。多项观察性研究发现，报告使用这些激素的女性心脏病发生率较低。护士健康研究在《新英格兰医学杂志》报告，雌激素联合孕激素使用者严重冠心病的调整后相对风险为0.39（95\%置信区间0.19～0.78）（17）。

\textbf{哪里出了问题：}研究比较了激素现用者与非使用者，这本身有问题，原因见“治疗或干预的队列研究”。此外，暴露时间判定还有一个错误：开始激素治疗后头一两年的心血管事件，被计入了未暴露女性！

这是因为激素治疗暴露根据参与者交回的问卷，每两年才更新一次。某女性1980年报告未用激素，1981年开始使用，在交回1982年问卷前，都被计为非使用者；到那时，若仍在使用，才开始累计使用者时间。若她在用激素第一年发生冠状动脉事件，即使是激素所致，也会被记在非使用者中！

这种暴露错分是有意为之。如表8.2所述，暴露测量先于结局，是增强前瞻性的特征。若女性发生冠状动脉事件后，才报告此前的激素使用，研究就不那么“前瞻”。作者写道：“为保持研究的前瞻性，每个两年时段的激素使用（包括持续时间），均根据该时段开始时女性的报告确定；因此，我们可能平均低估了约一年的使用时长。”（18）\sourcepage{228}

研究者直到随机试验提示后才认识到，激素似乎主要在头一两年增加冠状动脉事件。将NHS按一系列目标试验分析，比较开始激素治疗者与未开始者，这种早期不良效应才显现出来（19）。
\end{bookbox}

最后，如果研究的暴露（如他汀治疗）既由某协变量引起（如低密度脂蛋白〔low-density lipoprotein，LDL〕胆固醇高），又改变该协变量（他汀降低LDL胆固醇），队列研究的设计和解释就最为复杂。分析此类资料所需的高级方法，使多数研究者不愿采用，也使多数读者难以理解或信服结果报告。

\subsection{双队列、多队列研究与外部对照}
\term{双队列研究}（double-cohort study）有两个暴露水平不同的独立参与者队列。可从不同人群抽取（图8.5），更理想的是从同一人群抽取，此时称为\term{巢式}（nested）双队列研究。可选择等大队列；也可在某类暴露人数有限时，从另一人群（如未暴露者）抽取更多人，提高效能，例如未暴露与暴露人数比为2:1或3:1。（第6章曾用相似策略，相对于病例增加对照。）有时，按暴露队列中某些变量（如年龄和性别）的值匹配选取比较队列；这种匹配更常用于病例对照研究，见第9、10章。
\sourcepage{229}
\begin{figure}[htbp]\centering\chEightDouble
\caption*{\textbf{图8.5\quad 双队列研究。}从暴露（主要预测因素）水平不同、但符合相同纳入和排除标准的人群中选择两个队列；在基线测量其他预测因素；在结局发生时或随访结束时测量结局变量。}
\end{figure}

界定暴露和未暴露队列后，研究者测量其他协变量的基线水平，随访两个队列，在结局发生时、随访结束时，或暴露后特定间隔评估结局。因暴露水平已经界定随访组，这些研究通常简单或仅中等复杂，通常无须重复测量协变量。

双队列使用两个不同样本，不应与病例对照设计中的两个样本混淆（第9章）。双队列按暴露水平选择两组参与者；病例对照按结局有无选择。（遗憾的是，有人误将暴露组和未暴露组称为“病例”“对照”，并把研究误称为“病例对照研究”。）
\sourcepage{230}
\term{多队列研究}（multiple-cohort studies）与双队列相似，但超过两个队列，可以让同一暴露队列与两个以上未暴露队列比较，或让多个暴露队列与一个未暴露队列比较。例如，黄疸与婴儿喂养研究（Jaundice and Infant Feeding，JIFee）将一个未暴露队列（从产生暴露队列的出生队列中随机抽取）与两个暴露队列比较：一个胆红素高，另一个因脱水再次住院（例8.6）。

多队列设计的一种变体，是将一个或多个暴露队列的结局率，与从人口普查、登记或生命统计资料取得的\term{外部对照}（external controls）结局率比较。例如，多项研究考察男女医师自杀率，发现女医师自杀率高于一般女性，男医师则低于一般男性（23）。

\subsection{治疗或干预的队列研究}
无论研究暴露是危险因素还是疾病治疗，队列的基本结构相同；但差异仍足以单设本节，并另设一种工具评估治疗队列研究的偏倚风险（24）。

不同之一是，大多数疾病危险因素不能随机，而疾病治疗可以、通常也应该通过随机试验研究。队列研究，尤其回顾性队列，可能比随机试验更合乎伦理或更可行，因此仍可用于评估治疗，但通常应居于次要地位。

设计治疗队列时，应尽可能模拟为回答该问题而设计的随机试验（Hernán和Robins称之为“\term{目标试验}”，target trial）（24、25）。例如，这种试验的纳入和排除标准是什么？若有人因为随机分到某个治疗组将不合伦理而会被目标试验排除，那么他很可能也不应纳入该治疗的队列研究。随机试验随访起点通常清楚，即随机分组之日；复杂的治疗队列则不然。开始服药者可从首次配药起累计暴露随访时间，但未开始者的随访从何时起算？

如能找到让人符合目标试验资格的触发事件作为起点，会很有帮助。例如，癫痫发作后有人开始服药，有人没有，那么发作就是两组共同的随访起点；可以设想一个发作后才符合资格的目标试验。糖尿病患者第二次糖化血红蛋白A1c升高，或高血压患者第二次血压读数高，同样可以是触发事件，也可作为比较单药控制不佳者调整药物方案的目标试验纳入标准。

一旦确定目标试验资格的时点，就不应再用随后发生的事件改变已界定的暴露组，正如随机分到治疗的参与者必须按原随机组分析一样\sourcepage{231}（“一经随机，始终纳入分析”；见第11章）。例如，一项哮喘住院患者抗生素与住院时长的研究，最初允许暴露组在入院日（目标试验中的随机日）以后开始抗生素，且要求至少治疗两天。开始用药但不足两天便出院者，不计入抗生素组（26）。研究报告抗生素与住院中位时长增加一天有关。有人指出设计可能引入偏倚（27）后，研究者重新分析：只计第一天开始的抗生素，不再要求暴露组的最低治疗时长，住院中位时长差异随即消失（28）。

暴露组在接受两天抗生素前不可能出院，由此产生的偏倚称为\term{不死时间偏倚}（immortal time bias）（27、29、30）。这一名称源于最早报告此偏倚的生存研究：设计使暴露者在某些时段不可能死亡，却仍将这些时段计为其存活时间（31）。目标试验框架有助于识别和预防这种偏倚及“观察性研究中其他自己造成的伤害”（32）。

那么，研究开始时已经接受治疗（如已经服药）者如何处理？所谓\term{现用者}（prevalent users）的研究有数个问题（33—36）。第一，现用者中依从处方者占比过高；研究开始前已因任何原因停药者，不会被算作使用者。即使药物是安慰剂，且研究者控制了其他有利健康因素，依从处方者的结局仍较好（37、38）！这并非小效应：一项荟萃分析中，良好依从安慰剂者的合并死亡比值比为0.56（95\%置信区间0.43～0.74）（37）。因此，比较现用者与非使用者，会倾向于使任何药物——包括安慰剂——看起来有益。

第二，参与者可能因为许多不常测量的原因而更易接受治疗，如重视预防，或与医师关系密切。这些因素既影响处方，又影响结局风险，产生“\term{健康使用者}效应”（healthy user effect）。例如，队列研究中，获处方骨质疏松药物的女性，全因死亡看起来比可比未治疗者低25\%～60\%，但大型随机安慰剂对照试验未发现此益处（39）。

第三，纳入现用者会漏掉开始用药时的早期不良效应。若不良效应（包括死亡）早期发生，这一点尤其重要。例如，女性激素替代治疗的心血管疾病增加似乎主要在开始治疗后的第一年（19、40）。

第四，现用者研究使混杂控制更复杂（第10章），因为研究开始时混杂变量值可能已经被治疗改变。前述反复测量LDL胆固醇的他汀疗效研究就是例子；抗逆转录病毒现用者研究中控制病毒载量，或骨质疏松药物研究中控制骨密度，同样如此。\sourcepage{232}这些情况下，混杂变量基线值可能已受治疗影响，控制它反而会掩盖部分治疗益处。

最后，现用者与非使用者的比较无法对应到目标试验。对于尚未服药者，可以设想比较开始与不开始服药的目标试验，帮助决定是否开始；对已经服药者，可以设想比较继续与停药的试验，帮助决定是否停止。但不存在比较“已经在服药”与“没有在服药”之作用的目标试验。

\subsection{队列研究的优点与缺点}
队列设计的一大优点是，与横断面设计不同，它能计算发病，即随时间出现某状况的率或比例（表8.1）。纵向方法也可考察血压等连续结局的变化。在结局前测量预测因素，可确立时间顺序，加强从关联推断因果的基础。

所有队列都具有观察性研究相对于临床试验的一般弱点：因果推断困难，混杂变量常使解释含混（第10章）。前瞻性队列还有两项弱点：费用高、研究罕见结局效率低。即使肺癌等看似常见的疾病，每年的发病也很少，必须长期随访大量人群，才能积累足够结局，产生有意义的结果。对于较常见、较快发生的二分结局，以及连续结局，队列设计更有效率。

尽管治疗有效性和安全性研究会有多种偏倚，队列有时是估计某些疗效的唯一可行方法。例如，疾病罕见，或可选治疗极少甚至没有而人们坚信某治疗有效时，可能无法开展样本量充分的随机盲法试验。

有人认为，观察性研究纳入不符合临床试验资格者，可提供治疗效果与安全性更好的“真实世界”资料。这种研究可能发现临床试验未能发现的不良效应，尤其在合并其他疾病的老年患者中。但“真实世界”队列中有利的效果报告易受前述多种偏倚影响，必须审慎看待。

回顾性队列具有前瞻性队列的许多优点，且成本、耗时少得多：参与者已组建，基线测量已完成，随访期已经过去。主要缺点是，研究者对抽样与随访方法，以及测量性质和质量的控制有限；既有资料可能不完整、不准确，或其测量方式不适合研究问题。

\sourcepage{233}对于罕见职业或环境潜在危害暴露，双队列或多队列可能是唯一可行方法。用人口普查或登记作外部对照，还有基于人群且经济的好处；其他优点与一般队列类似。一个顾虑是，若从不同人群抽取暴露不同者，选择偏倚风险更高，因为除暴露外，队列还可能在影响结局的重要方面不同。年龄、种族等差异可匹配或统计调整，但地区等基础人群差异可能使两组不可比。因此，更理想的是暴露与未暴露都来自同一人群的巢式多队列，如阿奇霉素／阿莫西林研究（例8.5）和JIFee研究（例8.6）。若结局确定困难或昂贵而暴露罕见，巢式双队列或多队列可比研究整个队列高效得多。

\begin{bookbox}{例8.5\quad 巢式双队列研究}
人们担心常用门诊抗生素阿奇霉素（azithromycin）可能增加心源性猝死，而另一常用抗生素阿莫西林（amoxicillin）没有这种报告。Zaroff等为研究此风险（20）：
\begin{enumerate}
\item \textbf{确定抽取参与者的队列：}北加州或南加州Kaiser Permanente两个大型健康计划的参保患者，1998年1月1日至2014年12月31日有就诊和处方电子记录。
\item \textbf{界定暴露组：}找出处方前已连续参保且有处方福利至少12个月的患者中，门诊阿奇霉素处方（约170万张）或阿莫西林处方（约610万张）。排除处方前存在严重基础疾病者；另从电子病历取得多项人口学和临床潜在混杂变量（第10章）。
\item \textbf{界定并测量两组结局：}主要结局为处方后0～5天及6～10天的心血管死亡，来自死亡证明诊断编码；由心脏科医师组判定心血管死亡，以及随机抽取的非心血管死亡样本。
\end{enumerate}
阿奇霉素处方者5天内心血管死亡风险显著较高（调整后风险率比1.82；95\%置信区间1.23～2.67），6～10天则未见增加。但作者提醒，因果性不清楚，尤其是非心血管死亡也有相似关联（调整后风险率比2.17；95\%置信区间1.44～3.26）。

\textbf{点评：}回顾性设计允许从大人群16年的资料中捕获罕见事件，提供充分效能检验假设。暴露与结局也很可能测量准确。
\end{bookbox}

\sourcepage{234}
\begin{bookbox}{例8.6\quad （巢式）多队列设计}
为确定严重新生儿黄疸或脱水是否对神经发育有细微不良影响，UCSF和北加州Kaiser Permanente（KPNC）研究者开展三队列研究（21、22）。步骤如下：
\begin{enumerate}
\item \textbf{确定不同暴露的队列。}用电子数据库找出1995—1998年在KPNC医院出生的足月和近足月儿（$N=106627$），其中：
\begin{enumerate}
\item 最高血清总胆红素至少25 mg/dL（$N=147$）；或
\item 因脱水再次住院，血清钠至少150 mEq/L，或体重较出生时下降至少12\%（$N=197$）；或
\item 从出生队列随机抽取（$N=428$，未抽中暴露者）。
\end{enumerate}
\item \textbf{应用排除标准。}此步需要审查纸质病历并联系初级保健提供者，因此只能在队列抽样后完成。（排除了7名高胆红素血症者、15名脱水者和9名对照。）
\item \textbf{收集结局资料。}检索电子数据库中的神经系统疾病诊断；对同意者在5岁时开展全面神经发育检查，检查者不知道儿童属于哪个队列。
\end{enumerate}
高胆红素血症与脱水均未见与细微不良结局相关。
\end{bookbox}

针对暴露队列（41）和干预队列（24），已有评估偏倚风险的工具。计划队列研究者可查阅，确保至少在设计阶段处理容易补救的效度威胁。

\subsection{队列研究的分析方法}
\subsubsection{结局频率指标}
\term{风险}、\term{比值}与\term{率}，都是随访一定时间后，参与者二分结局频率的估计。三者密切相关，分子相同，均为发生结局的人数；都隐含\term{处于风险中}（at risk）的概念，即研究开始时尚无所关注的结局。前瞻性心绞痛预测因素研究中，基线已有心绞痛者不处于风险中，因为他已经有结局。观察到的风险，是发生结局人数除以风险人数；观察到的比值，是发生结局人数除以未发生结局人数，也等于风险／（1－风险）。

\sourcepage{235}许多队列中会出现失访、死亡，或其他使结局无法确定的事件。研究者须测量每人贡献的人时：从入队列开始，直到发生关注结局，或因失访、死亡而删失。某组（如暴露组）的发病密度，是该组结局数除以风险人时之和。

设1000人随访两年，观察肺癌发生，每年有50例新病例。表8.4列出风险、比值和率。

\begin{table}[htbp]\centering\small
\caption*{\textbf{表8.4\quad 1000人随访两年、每年50例新病例的风险、比值和率的简化计算}}
\begin{tabularx}{\textwidth}{L{18mm}Y Y}\toprule
统计量&公式&计算\\\midrule
风险&发生结局人数／风险人数&$100/1000=0.10$（10\%）\\
比值&发生结局人数／未发生结局人数&$100/900=0.111$\\
率&发生结局人数／风险人时$^{a}$&$100/1900$人年＝每年5.3\%，或每100人年5.3例\\\bottomrule
\end{tabularx}
\par\smallskip\raggedright\footnotesize
$^{a}$分母约为第一年的风险人年（年初1000人、年末950人，约975人年），加第二年（年初950人、年末900人，约925人年），共1900人年。此例每年病例数不变，而风险人数减少，因此率随时间略升。
\end{table}

风险最容易理解，因为日常很熟悉：两年肺癌风险为100／1000（10\%）。比值较不直观，为100:900，即0.111。罕见结局的比值与风险数值接近；此例10\%已接近二者仍可视为相似的上限。率考虑了第一年已经发生结局者，第二年不再处于风险中，不能列入第二年分母。因此第二年结局率（50例／约925人年）略高于第一年。

同一人可能多次发生某结局，例如反复链球菌性咽炎，或反映疾病加重的结局，如心力衰竭住院。若多数人只有0或1次，将其二分为0次与至少1次，信息损失很小而更简单。若许多人反复发生，则可计算每人平均结局次数，或将平均次数除以风险人时。\footnote{若有些人发生多次结局，不能简单混合不同人的人时，因为同一人的多次结局并不独立。若这句话涉及你不熟悉的内容，宜请教统计学家。}

\subsubsection{关联指标}
两组比较中，\term{风险差}（risk difference）就是风险的差，通常为暴露组减未暴露组。若暴露预防而非引起结局（例如某治疗），风险差为负，其绝对值称为\term{绝对风险降低}（absolute risk reduction）。

\sourcepage{236}两组风险的商，通常为暴露除以未暴露，称为\term{风险比}或\term{相对风险}（risk ratio／relative risk，RR）。\term{比值比}为暴露组结局比值与未暴露组比值之比，结局罕见时类似风险比。但比值容易令人困惑，尤其结局常见时，所以队列分析最好避免比值比。例外是使用\term{logistic回归}（logistic regression），这种多变量方法本身用比值比量化关联。\footnote{横断面和病例对照研究（第9章）常使用比值比；这时它可能是患病比值之比，而非发病比值之比。}

与风险比相似，\term{率比}（rate ratio）可估计为有、无某危险因素者的率之商。\term{Cox比例风险模型}（Cox proportional hazard model）可对这种有时称为“事件发生时间”的资料进行多变量分析，估计\term{风险率比}（hazard ratios）；它类似率比，是队列研究常用的关联指标。

\subsection{尽可能完成随访}
随访整个队列很重要，前瞻性研究应采取数项措施（表8.5）。研究者、工作人员和参与者关系良好，且访视体验积极、愉快、有参与感，是完整随访的重要条件；参与者往往期待与工作人员或研究联系。含人物介绍、食谱或其他消息的通讯，可增强对研究的归属感。

\begin{bookbox}{表8.5\quad 尽量减少失访的策略}
\textbf{入组期间}
\begin{enumerate}
\item 建立工作人员与参与者之间的关系：使基线和随访访视有趣、愉快、有参与感。
\item 排除可能失访者：计划搬迁；不能确定是否愿意回来；身体欠佳或有与研究问题无关的致命性疾病。
\item 获取便于日后寻找的信息：参与者的地址、电话（手机号码尤有用）和电子邮箱；社会保障／Medicare号码，以及未来追踪结局时使用的许可；不与其同住的密友或亲属姓名、地址、电话和邮箱；医师姓名、地址、电话和邮箱。
\end{enumerate}
\textbf{随访期间}
\begin{enumerate}
\item 定期联系，收集资料、提供结果并给予支持：电话可能需在周末、晚间打；电子邮件或附已付邮资、写好回邮地址的卡片可反复寄送；通讯、生日卡和小礼物；若不影响结局，可告知随访检查结果。
\item 对电话和信件联系不到者：联系朋友、亲属或医师；向邮政部门查询转寄地址；\sourcepage{237}通过电话簿、互联网等公开来源找地址，若最初获得许可，也可查询征信机构；对Medicare参与者，从社会保障署收集出院资料；从州卫生部门或全国死亡索引确定生死状况。
\end{enumerate}
\textbf{任何时候}

感谢参与者，友善且尊重地对待他们，帮助他们理解研究问题，使他们愿意以伙伴身份共同促成研究成功。
\end{bookbox}

需要面对面测量的研究，应在一开始排除研究期间计划搬迁至无法联系地区者。尽早收集便于找到迁居者或死者的信息，包括参与者、其私人医师及不同住的亲友地址、电话和邮箱。手机号码和个人邮箱特别有用，因为搬家或换工作时常不变。如可行，取得社会保障号码有助于确定失访者生死状况，并取得Medicare参与者的社会保障署出院资料。至少每年联系一两次，有助于掌握去向，提高结局记录的及时性和准确度。有时须坚持反复发信、发邮件、打电话，甚至家访，才能完成随访评估。

即使尽力，多数研究随访仍不足100\%。此时，\term{敏感性分析}（sensitivity analyses）可考察结论是否取决于对失访者结局的假定，例如假设暴露组失访者均未发生、未暴露组失访者全发生结局，或反过来，重新分析。若结论不随这些假定改变，研究就更有说服力。

\section{小结}
\begin{enumerate}
\item 横断面研究在同一时点测量全部变量，预测因素与结局没有结构性区分。因不能证明预测变量先于结局，其因果依据弱于队列研究。
\item 横断面研究能提供有价值的患病率描述，避免随访设计的时间、费用和退出问题；可作为队列或实验的第一步，也可将独立抽样的连续调查连接，以显示人群随时间的变化。
\item 队列研究随访一开始确定的参与者，描述某状况的发病或自然史，寻找各种结局的预测因素（危险因素）。
\item 队列通常分前瞻性与回顾性，但术语使用不一；应明确暴露与结局如何、何时测量，才能充分理解偏倚风险。
\sourcepage{238}
\item 前瞻性队列一般从随访起点开始，可能须长期招募和随访大量参与者；测量控制更好，但常耗时且昂贵。
\item 如能找到预测变量与结局都已测量的回顾性队列，有时可克服此缺点，但测量控制较少。
\item 双队列或多队列比较暴露性质或水平不同队列的结局发生，适合罕见暴露与职业暴露。巢式回顾性多队列适合研究同一人群内罕见暴露的作用。
\item 参照可行时本应开展的治疗随机试验，即目标试验，构建治疗或干预队列，可减少偏倚风险。
\item 风险、比值和率都能估计随访中二分结局的频率；考虑风险人时的发病密度，是使用Cox比例风险模型计算多变量风险率比的现代方法的基础。
\item 随访不全会削弱队列优势；一开始排除可能无法随访者，收集便于寻找的基线资料，并定期联系每个人，可减少失访。
\end{enumerate}

\begingroup\sloppy
\begin{references}
\item Drouin-Chartier JP, Chen S, Li Y, et al. Egg consumption and risk of cardiovascular disease: three large prospective US cohort studies, systematic review, and updated meta-analysis. BMJ. 2020;368:m513.
\item Centers for Disease Control NCHS. National Health and Nutrition Examination Survey (NHANES): History 2020 [cited 2020 Dec 15]. \url{https://www.cdc.gov/nchs/nhanes/history.htm}.
\item Andersen RE, Crespo CJ, Bartlett SJ, Cheskin LJ, Pratt M. Relationship of physical activity and television watching with body weight and level of fatness among children: results from the Third National Health and Nutrition Examination Survey. JAMA. 1998;279(12):938-942.
\item Sargent JD, Beach ML, Adachi-Mejia AM, et al. Exposure to movie smoking: its relation to smoking initiation among US adolescents. Pediatrics. 2005;116(5):1183-1191.
\item Morgenstern M, Sargent JD, Engels R, et al. Smoking in movies and adolescent smoking initiation: longitudinal study in six European countries. Am J Prev Med. 2013;44(4):339-344.
\item Dal Cin S, Stoolmiller M, Sargent JD. Exposure to smoking in movies and smoking initiation among black youth. Am J Prev Med. 2013;44(4):345-350.
\item Mejia R, Perez A, Pena L, et al. Smoking in movies and adolescent smoking initiation: a longitudinal study among Argentinian adolescents. J Pediatr. 2017;180:222-228.
\item Blackwell T, Yaffe K, Ancoli-Israel S, et al. Association of sleep characteristics and cognition in older community-dwelling men: the MrOS sleep study. Sleep. 2011;34(10):1347-1356.
\item Kalantar-Zadeh K, Abbott KC, Salahudeen AK, Kilpatrick RD, Horwich TB. Survival advantages of obesity in dialysis patients. Am J Clin Nutr. 2005;81(3):543-554.
\item McMillen RC, Gottlieb MA, Shaefer RM, Winickoff JP, Klein JD. Trends in electronic cigarette use among U.S. adults: use is increasing in both smokers and nonsmokers. Nicotine Tob Res. 2015;17(10):1195-1202.
\item Vandenbroucke JP. Prospective or retrospective: what’s in a name? BMJ. 1991;302(6771):249-250.
\item Rothman KJ, Greenland S, Lash TL. Modern Epidemiology. 3rd ed. Wolters Kluwer Health/Lippincott Williams \& Wilkins; 2008:758.
\sourcepage{239}
\item Bao Y, Bertoia ML, Lenart EB, et al. Origin, methods, and evolution of the three nurses’ health studies. Am J Public Health. 2016;106(9):1573-1581.
\item Man S, Xian Y, Holmes DN, et al. Association between thrombolytic door-to-needle time and 1-year mortality and readmission in patients with acute ischemic stroke. JAMA. 2020;323(21):2170-2184.
\item Browner WS, Pressman AR, Nevitt MC, Cauley JA, Cummings SR. Association between low bone density and stroke in elderly women: the study of osteoporotic fractures. Stroke. 1993;24(7):940-946.
\item Evans DS, Kim KM, Jiang X, Jacobson J, Browner W, Cummings SR. Prediction of in-hospital mortality among adults with COVID-19 infection. medRxiv. 2021.
\item Grodstein F, Stampfer MJ, Manson JE, et al. Postmenopausal estrogen and progestin use and the risk of cardiovascular disease. N Engl J Med. 1996;335(7):453-461.
\item Grodstein F, Manson JE, Colditz GA, Willett WC, Speizer FE, Stampfer MJ. A prospective, observational study of postmenopausal hormone therapy and primary prevention of cardiovascular disease. Ann Intern Med. 2000;133(12):933-941.
\item Hernán MA, Alonso A, Logan R, et al. Observational studies analyzed like randomized experiments: an application to postmenopausal hormone therapy and coronary heart disease. Epidemiology. 2008;19(6):766-779.
\item Zaroff JG, Cheetham TC, Palmetto N, et al. Association of azithromycin use with cardiovascular mortality. JAMA Netw Open. 2020;3(6):e208199.
\item Newman TB, Liljestrand P, Jeremy RJ, et al. Outcomes among newborns with total serum bilirubin levels of 25 mg per deciliter or more. N Engl J Med. 2006;354(18):1889-1900.
\item Escobar GJ, Liljestrand P, Hudes ES, et al. Five-year neurodevelopmental outcome of neonatal dehydration. J Pediatr. 2007;151(2):127-133, 133 e1.
\item Duarte D, El-Hagrassy MM, Couto TCE, Gurgel W, Fregni F, Correa H. Male and female physician suicidality: a systematic review and meta-analysis. JAMA Psychiatry. 2020;77(6):587-597.
\item Sterne JA, Hernán MA, Reeves BC, et al. ROBINS-I: a tool for assessing risk of bias in non-randomised studies of interventions. BMJ. 2016;355:i4919.
\item Hernán MA, Robins JM. Using big data to emulate a target trial when a randomized trial is not available. Am J Epidemiol. 2016;183(8):758-764.
\item Stefan MS, Shieh MS, Spitzer KA, et al. Association of antibiotic treatment with outcomes in patients hospitalized for an asthma exacerbation treated with systemic corticosteroids. JAMA Intern Med. 2019;179(3):333-339 [original version, subsequently retracted].
\item Newman TB. Possible immortal time bias in study of antibiotic treatment and outcomes in patients hospitalized for asthma. JAMA Intern Med. 2021;181(4):568-569.
\item Stefan MS, Shieh M-S, Spitzer KA, et al. Association of antibiotic treatment with outcomes in patients hospitalized for an asthma exacerbation treated with systemic corticosteroids. JAMA Intern Med. 2019;179(3):333-339 [replaced version of retracted article].
\item Newman T, Kohn M. Evidence-Based Diagnosis: An Introduction to Clinical Epidemiology. 2nd ed. Cambridge University Press; 2020:231-243.
\item Newman TB. Antibiotic treatment for inpatient asthma exacerbations: what have we learned? JAMA Intern Med. 2021;181(4):427-428.
\item Suissa S. Immortal time bias in pharmaco-epidemiology. Am J Epidemiol. 2008;167(4):492-499.
\item Hernán MA, Sauer BC, Hernandez-Diaz S, Platt R, Shrier I. Specifying a target trial prevents immortal time bias and other self-inflicted injuries in observational analyses. J Clin Epidemiol. 2016;79:70-75.
\item Vandenbroucke J, Pearce N. Point: incident exposures, prevalent exposures, and causal inference: does limiting studies to persons who are followed from first exposure onward damage epidemiology? Am J Epidemiol. 2015;182(10):826-833.
\item Hernán MA. Counterpoint: epidemiology to guide decision-making: moving away from practice-free research. Am J Epidemiol. 2015;182(10):834-839.
\item Vandenbroucke J, Pearce N. Vandenbroucke and Pearce respond to “incident and prevalent exposures and causal inference”. Am J Epidemiol. 2015;182(10):846-847.
\sourcepage{240}
\item Ray WA. Evaluating medication effects outside of clinical trials: new-user designs. Am J Epidemiol. 2003;158(9):915-920.
\item Simpson SH, Eurich DT, Majumdar SR, et al. A meta-analysis of the association between adherence to drug therapy and mortality. BMJ. 2006;333(7557):15.
\item Avins AL, Pressman A, Ackerson L, Rudd P, Neuhaus J, Vittinghoff E. Placebo adherence and its association with morbidity and mortality in the studies of left ventricular dysfunction. J Gen Intern Med. 2010;25(12):1275-1281.
\item Cummings SR, Lui LY, Eastell R, Allen IE. Association between drug treatments for patients with osteoporosis and overall mortality rates: a meta-analysis. JAMA Intern Med. 2019;179(11):1491-1500.
\item Hulley S, Grady D, Bush T, et al. Randomized trial of estrogen plus progestin for secondary prevention of coronary heart disease in postmenopausal women. Heart and Estrogen/progestin Replacement Study (HERS) Research Group. JAMA. 1998;280(7):605-613.
\item Bero L, Chartres N, Diong J, et al. The risk of bias in observational studies of exposures (ROBINS-E) tool: concerns arising from application to observational studies of exposures. Syst Rev. 2018;7(1):242.
\end{references}
\endgroup

\clearpage
\chapterappendix{8A}{第8章习题：横断面研究与队列研究的设计}
\sourcepage{241}
\begin{enumerate}
\item 研究问题为：“维生素B12水平低会导致老年人髋部骨折吗？”
\begin{enumerate}
\item 简述如何用前瞻性队列设计回答。
\item 另一方法是从老年门诊抽样，比较既往有髋部骨折者与所有未骨折者的维生素B12。相对于此横断面方法，列出前瞻性队列至少一项优点和一项缺点。
\item 能否用回顾性队列回答？若能，对上述优缺点有何影响？
\end{enumerate}
\item Sung等（1）研究338名至少30岁、超重或肥胖女性在首次入组饮食和运动减少尿失禁项目（Program to Reduce Incontinence by Diet and Exercise，PRIDE）临床试验时，尿失禁频率与抑郁症状的关联。有抑郁症状者（$N=101$）报告每周失禁次数多于无症状者（28次比23次；$P=0.005$）。
\begin{enumerate}
\item 这是什么研究设计？
\item 一种解释是抑郁增加尿失禁频率。还可能有哪些解释？如何改变设计来区分？
\end{enumerate}
\end{enumerate}
\begin{references}
\item Sung VW, West DS, Hernandez AL. Association between urinary incontinence and depressive symptoms in overweight and obese women. Am J Obstet Gynecol. 2009;200(5):557.e1-557.e5.
\end{references}
